%%%%%%%%%%%%%%%%%%%%%---Chapter1 Formatting---%%%%%%%%%%%%%%%%%%%%%%%%%%%%%%%%%%
\documentclass[Thesis.tex]{subfiles}

\begin{document}
\newpage
\setcounter{chapter}{0} %%Chapter Numbering

% %%Adjust chapter title formatting
\titleformat{\chapter}[display]
{\fontsize{27}{27}\bfseries\filleft}{ \thechapter}{0pt}{{#1}}

\thispagestyle{empty}
%%%%%%%%%%%%%%%%%%%%%%%%%%%%%%%%%%%%%%%%%%%%%%%%%%%%%%%%%%%%%%%%%%%%%%%%%%%

\chapter{Introduction}\label{chap1:introduction}

\section{Introduction}

Quantum computers available today are noisy. Every gate application, every idle
period and every measurement introduces errors at rates that, while remarkable by
the standards of a decade ago, remain several orders of magnitude above what
fault-tolerant quantum computation requires. A circuit of even modest depth
accumulates enough error that the quantity a user actually wants -- an expectation
value $\langle O \rangle$ of some observable on the ideal output state -- is
displaced far from its noiseless value. This regime, in which devices are large
enough to be interesting but too noisy to be trusted, is what the literature calls
the noisy intermediate-scale quantum (NISQ) era~\cite{preskill2018nisq,bharti2022noisy}, and it is the setting for the whole
of this thesis.

Two distinct research programmes address this problem, and they differ in what they
promise. \emph{Quantum error correction} (QEC) encodes one logical qubit into many
physical qubits and, provided the physical error rate is below a code-dependent
threshold, suppresses the logical error rate exponentially in the code
distance~\cite{dennis2002topological,fowler2012surface,google2025below}. It
is the asymptotically correct answer, and it is expensive: the qubit overhead is
large and the classical decoding problem must be solved in real time, within the
duration of a syndrome-extraction round. \emph{Quantum error mitigation} (QEM) makes
no attempt to correct the quantum state. Instead it accepts the noisy device as
given and post-processes the measured data so that the \emph{estimator} of an
expectation value is less biased than the raw measurement~\cite{temme2017error,li2017efficient,cai2023review}.
Mitigation needs no extra qubits, which is why it is the technique of choice on present hardware, and it buys
that cheapness with sampling overhead rather than qubit overhead -- an overhead
that is known to grow exponentially with circuit
depth~\cite{takagi2022fundamental,quek2024exponentially}.

This thesis is a benchmarking study of machine-learning approaches to both problems,
with particular attention to \emph{hybrid quantum--classical} models -- models in
which a parametrized quantum circuit does part of the computation and a classical
optimizer or regressor does the rest. The question it asks is deliberately narrow
and, as the results will show, not favourable to the hypothesis that motivated it.

\section{Error mitigation as an estimation problem}

The cleanest way to state the mitigation problem is as one of estimation. Let
$\rho_{\text{ideal}}$ be the state a circuit $C$ would produce on a perfect device
and $\rho_{\text{noisy}}$ the state it actually produces. For an observable $O$ we
can measure
\begin{equation}
    \langle O \rangle_{\text{noisy}} = \mathrm{Tr}\!\left[ O \rho_{\text{noisy}} \right],
    \label{eq:ch1_noisy}
\end{equation}
but we want
\begin{equation}
    \langle O \rangle_{\text{ideal}} = \mathrm{Tr}\!\left[ O \rho_{\text{ideal}} \right].
    \label{eq:ch1_ideal}
\end{equation}
Mitigation constructs an estimator $\widehat{\langle O \rangle}$ from measurements
taken on the noisy device -- possibly at several deliberately different noise levels,
possibly on several related circuits -- whose bias with respect to
Eq.~\eqref{eq:ch1_ideal} is smaller than that of Eq.~\eqref{eq:ch1_noisy}. Zero-noise
extrapolation, the most widely deployed technique, measures the observable at an
amplified noise level and extrapolates the resulting curve back to the zero-noise
axis~\cite{temme2017error,giurgicatiron2020digital,he2020zero}. It requires no training data and makes only a smoothness assumption about how
the expectation value depends on noise strength.

Once the problem is written this way, the machine-learning formulation is immediate.
If we can obtain, for a set of \emph{reference} circuits, both the noisy measurements
and the true ideal values, then the map from one to the other is a supervised
regression problem, and any regressor may be trained on it. The reference values can
come from classical simulation of circuits small enough to simulate, from
Clifford-dominated circuits that are efficiently simulable by the Gottesman--Knill
theorem~\cite{nielsen2012book}, or from circuits whose observable is known
analytically. This is the strategy behind Clifford data regression and its
learning-based descendants~\cite{czarnik2021clifford,strikis2021learning,lowe2021unified,liao2024machine}, and it
is the strategy this thesis benchmarks.

The appeal of a \emph{quantum} model for this regression is that the function being
learned is generated by quantum dynamics, so a model whose inductive bias is itself
quantum might represent it more naturally than a generic classical
regressor~\cite{havlicek2019supervised,schuld2021kernel,huang2021power}. That
intuition is plausible, widely repeated~\cite{schuld2022isquantum,bowles2024better}, and -- as Chapter~\ref{chap5:results} will
show -- not supported by the measurements reported here.

\section{Research motivation}

The literature on machine-learning error mitigation reports encouraging results, but
three methodological weaknesses recur often enough that the aggregate picture is
difficult to trust.

\textbf{Baselines are inconsistent.} Papers report improvement over the unmitigated
expectation value, which is a low bar, far more often than they report improvement
over a properly tuned zero-noise extrapolation, which is the technique a
practitioner would otherwise use. A learned model that beats doing nothing has not
yet earned its place in a workflow.

\textbf{Cost is rarely normalized.} A mitigation method's price is quantum: how many
circuit evaluations must be executed per estimated expectation value. Zero-noise
extrapolation with three noise settings costs three evaluations. A learned model
consuming the same three noise-scaled inputs costs the same three, and is therefore
directly comparable; a model that additionally requires its own quantum circuits
costs more -- one extra circuit per prediction for a variational model, and one per
\emph{training point} for a fidelity-kernel model (Section~\ref{sec:ch5_cost}) -- and
is not. Reporting accuracy without this denominator makes
methods look comparable when they are not.

\textbf{Generalization is rarely tested.} A mitigation model is trained on circuits
that can be simulated or otherwise referenced, and then deployed on circuits that
cannot -- which, in practice, means deeper circuits, different circuit families, or
a device whose calibration has drifted since the training set was collected. The
in-distribution test that most papers report says nothing about that deployment
scenario, and the generalization theory of quantum models offers no guarantee
outside the training distribution~\cite{caro2022generalization,caro2023ood}. This thesis treats out-of-distribution evaluation as a first-class
experiment rather than an afterthought, and finds that it reverses the ranking of
methods.

A fourth motivation is specific to quantum kernel methods. Their performance depends
critically on a bandwidth hyperparameter that controls the spread of the embedded
data, and an untuned quantum kernel
concentrates~\cite{shaydulin2022importance,canatar2023bandwidth,thanasilp2024exponential}:
all off-diagonal entries of the Gram matrix collapse towards a constant, the model loses its ability to
discriminate, and the resulting regressor can be worse than no model at all. A
baseline reported without bandwidth tuning is a broken baseline, and this thesis
quantifies exactly how broken.

\section{Research objectives}

The objectives of this work are:

\begin{enumerate}
    \item To formulate error mitigation as a supervised regression task with an
    evaluation protocol that does not leak information between training and test
    sets, and to construct a benchmark dataset with an exactly specified noise model
    and known ground truth.

    \item To benchmark, on identical folds, three families of mitigation method: the
    unmitigated baseline and analytic zero-noise extrapolation variants that require
    no training; classical machine-learning regressors; and hybrid
    quantum--classical models, specifically quantum fidelity-kernel ridge regression
    and a data-reuploading variational quantum circuit.

    \item To report accuracy together with cost, measuring cost in circuit
    evaluations per prediction as well as in trainable parameters and training time,
    so that no method appears to win by spending quantum resources that its
    competitors were not allowed.

    \item To measure out-of-distribution behaviour explicitly, by testing on circuits
    deeper than any seen in training and on a circuit family not seen in training,
    and to determine which model classes degrade gracefully and which fail.

    \item To characterize the sensitivity of the hybrid models to their
    hyperparameters and design choices -- in particular quantum-kernel bandwidth and
    the periodicity of the variational encoding -- and to the training split and
    initialisation, to the point where an observed failure can be attributed to a
    cause rather than merely reported.

    \item To extend the study to quantum error correction by benchmarking
    machine-learning syndrome decoders, including a hybrid quantum--classical
    decoder, against minimum-weight perfect matching on syndrome
    data~\cite{edmonds1965paths,higgott2022pymatching,torlai2017neural,bausch2024learning},
using logical error rate as the metric.
\end{enumerate}

\section{Research contributions}

\subsection{Summary of contributions}

The substantive contributions of this thesis are as follows.

\begin{enumerate}
    \item \textbf{A two-regime benchmark design.} The main comparison, the
    noise-scaling comparison and the robustness analysis are run both under an
    idealized gate-local Markovian model -- the regime in which the machine-learning
    mitigation literature reports its results -- and under a device-realistic model
    with per-qubit inhomogeneous rates, coherent over-rotation, always-on $ZZ$
    crosstalk and non-Markovian dynamics. The bandwidth study and the five sweeps are
    run in the idealized regime only. A conclusion that survives only one regime is a
    conclusion about the noise model, and the design makes that visible.

    \item \textbf{Non-Markovianity that is measured and isolated.} Memory is
    introduced through a coherently coupled two-level-system defect and quantified
    with the Breuer--Laine--Piilo measure on an idling register. Two Markovian
    controls return exactly zero: one removes the defect and one resets it every
    layer. Only the reset control holds the defect's per-layer loss fixed, and using
    it reverses a conclusion the defect-removed control alone would have suggested:
    memory narrows the learned advantage over extrapolation rather than widening it.

    \item \textbf{A calibration-drift experiment}, in which mitigators are trained
    against one device calibration and evaluated after a $30\%$ recalibration on the
    same circuits. Most learned arms lose no resolvable accuracy; the raw error,
    linear extrapolation and ridge regression do.

    \item \textbf{Both an idealized and a hardware-realizable noise-scaling method.}
    Parameter scaling is exact but no device can perform it; layer-level unitary
    folding is what practitioners use. Under folding the learned arms' margin over
    linear extrapolation roughly halves but remains resolved in every split for the
    four non-variational learned arms.

    \item \textbf{Shot-resolved cost accounting for the quantum kernel.} The fidelity
    Gram matrix is re-estimated from finite shots as hardware must. Approaching the
    simulator accuracy needs about $10^5$ shots per entry, or $3 \times 10^8$ circuit
    executions per mitigated value.

    \item \textbf{A leakage-free protocol with paired statistics.} Every split is
    grouped by circuit, because the observables measured on one circuit share its
    noise realization. Methods are compared by a paired, circuit-grouped bootstrap on
    the difference of their errors, every headline comparison is repeated over ten
    independent splits, and the variational arms over ten initialisation seeds.

    \item \textbf{Learned mitigation removes $41$--$44\%$ of the residual error left
    by linear extrapolation in-distribution} in the idealized regime, resolved in
    every split, at the same circuit cost for the classical regressors.

    \item \textbf{A negative result on quantum advantage in this task.} The quantum
    fidelity-kernel regressor is tied with the classical neural network in the
    idealized regime and under folding, and ahead of it by about $3\%$ only under
    idealized parameter scaling in the realistic regime, while requiring one
    shot-estimated circuit per training point for every prediction. The
    data-reuploading variational circuit is unstable across initialisation seeds and
    on average no better than no mitigation in the idealized regime.

    \item \textbf{A quantification of quantum-kernel bandwidth sensitivity}: the
    untuned kernel is worse than no mitigation, and tuning bandwidth changes its error
    by a factor of $3.6$.

    \item \textbf{An out-of-distribution failure analysis with an identified cause.}
    On deeper circuits the neural network, the quantum kernel and the variational
    circuits are worse than no mitigation when circuit depth and gate count are among
    their inputs. Removing those two inputs, whose test values lie outside the
    training range, makes every learned arm resolvably better than no mitigation.

    \item \textbf{A syndrome-decoding benchmark} comparing neural and hybrid
    quantum--classical decoders against minimum-weight perfect matching on the
    distance-$3$ to $9$ repetition code under phenomenological noise, reported in
    logical error rate with decoding latency alongside. The repetition code
    substitutes for the surface code, for reasons stated in
    Chapter~\ref{chap6:decoding}.
\end{enumerate}

\subsection{Organization of thesis}

Chapter~\ref{chap2:background} reviews the noise models, mitigation techniques,
quantum machine-learning primitives and decoding algorithms on which this work
builds, and states the research gap precisely.
Chapter~\ref{chap3:methodology} formalizes the benchmark: the noise model, the
circuit families, the feature construction, the evaluation protocol, every model
arm, and the metrics.
Chapter~\ref{chap4:setup} documents the implementation -- software versions, seeds,
hyperparameter selection protocol and computational budget -- in enough detail to
reproduce every number in the thesis.
Chapter~\ref{chap5:results} reports the mitigation benchmark: the main comparison,
the noise-strength, shot-budget and system-size sweeps, learning curves, feature
ablation and the two out-of-distribution experiments in the idealized regime, and
then repeats the main comparison under a device-realistic, non-Markovian noise model
with hardware-realizable noise scaling, calibration drift and a shot-limited
quantum kernel.
Chapter~\ref{chap6:decoding} reports the syndrome-decoding track.
Chapter~\ref{chap7:conclusion} states the conclusions, the limits of what has been
established, and the directions that follow from the failure modes identified here.
Appendix~\ref{AppendixA} collects the derivations,
Appendix~\ref{AppendixB} the complete result tables and selected hyperparameters,
and Appendix~\ref{AppendixC} the principal source listings.

\end{document}
